% ============================================================
\section{Conclusion}
\label{sec:conclusion}
% ============================================================

This paper presented GNNocRoute-DRL, a framework combining Graph Neural Networks with Deep Reinforcement Learning for topology-aware adaptive routing in Network-on-Chip. Our contributions span three areas:

\textbf{First}, we demonstrated through graph-theoretic analysis across seven topologies that deterministic routing creates structurally predictable congestion patterns, with congestion imbalance reaching 0.475 on an $8\times8$ mesh. \textbf{Second}, we showed that GATv2-based GNN encoders learn topology-aware node embeddings that strongly correlate with betweenness centrality ($|r| = 0.978$), providing the DRL agent with structural information that MLP-based encoders cannot capture. \textbf{Third}, through extensive BookSim2 simulations spanning 252 configurations with power analysis, we demonstrated that adaptive routing reduces hotspot latency by up to 5.4\% compared to XY routing with negligible power overhead, while cross-topology analysis confirms GNN embedding transferability (Pearson $r \approx 0.96$).

The periodic optimization framework addresses the fundamental tension between GNN expressiveness ($10^6$ cycles per inference) and NoC timing constraints ($<10$ cycles per routing decision), achieving practical feasibility with $P = 5,000$ cycle update periods. Cross-topology transfer learning further demonstrates that GNN embeddings generalize across mesh sizes (Pearson $r \approx 0.96$).

Our results establish that GNN-enhanced topology awareness is both feasible and beneficial for NoC routing, providing a foundation for next-generation adaptive routing in billion-transistor chips.

\subsection*{Future Work}
Several directions warrant further investigation: (1) full DRL agent implementation and evaluation on BookSim2; (2) energy-aware routing with power model integration; (3) evaluation on real PARSEC benchmark workloads; and (4) hardware implementation of the GNN encoder for on-chip deployment.

\subsection*{Acknowledgments}
This research was supported by the Institute of Information Technology, Vietnam National University, Hanoi, under PhD research program No. 25218003.
